
% Sección 1: Introducción
\section{Introducción}
\begin{frame}{Contexto}
    \begin{itemize}
        \item \alerta{Auge de LLMs} (ChatGPT) en redacción científica.
        \item \textbf{Problemas éticos:}
              \begin{itemize}
                  \item Generación de datos, gráficos, código o texto que argumente, o sustente una idea y/o hipótesis no sustentada con datos reales.
                  \item Falta de transparencia en contenido generado con IA.
                  \item Alucinaciones (datos/citas falsas, no intencionales).
                  \item Artículos 100\% IA con ediciones mínimas.
              \end{itemize}
        \item \alerta{Artículo clave:} (\citeauthor{Kendall2024}, \citeyear{Kendall2024}) propone incluir \textbf{prompts y texto generado}\cite{Kendall2024}.
    \end{itemize}
    \begin{tikzpicture}
        \node[aLabel, text width=0.90\textwidth, align=left] (origin)  at (0,0){\small \fullcite{Kendall2024}};
    \end{tikzpicture}
\end{frame}
\begin{frame}{Introducción}
    \begin{itemize}
        \item Las revistas científicas requieren reconocer el uso de Modelos de Lenguaje Grande (LLMs).
        \item El problema: \textbf{las directrices actuales son insuficientes} \cite{Kendall2024}.
        \item Los revisores y lectores no saben:
              \begin{itemize}
                  \item Qué partes\footnote{Secciones,texto,imágenes, código, citas, etc.} fueron generadas con IA.
                  \item Cómo se editó el texto generado.
              \end{itemize}
        \item Este artículo (\citeauthor{Kendall2024}, \citeyear{Kendall2024}) propone cambios significativos en las directrices actuales.
    \end{itemize}
\end{frame}

% El problema de las directrices actuales
\begin{frame}{El problema de las directrices actuales}
    \begin{itemize}
        \item Permiten reconocer el uso de IA sin especificar:
              \begin{itemize}
                  \item Qué partes del artículo fueron generadas con IA.
                  \item Cuánto contenido fue asistido por IA (¿1\% o 100\%?).
                  \item Qué ediciones se realizaron al texto generado.
                  \item Qué herramienta de IA se ha utilizado, y en cuál de sus versiones.
              \end{itemize}
    \end{itemize}
    \begin{boxA}{Ejemplo}
        Hay artículos generados parcial o completamente por IA que cumplen con las directrices actuales de algunas revistas.\\[1em]
        {\small Como veremos en los casos de los autores M. Majovský y S. Biswas.}
    \end{boxA}
\end{frame}

% La cuestión de la autoría
\begin{frame}{La cuestión de la autoría}
    \begin{itemize}
        \item \textbf{¿Hay un límite después del cual no se puede reclamar autoría?}
              \begin{itemize}
                  \item COPE, ICMJE no lo contemplan. La mayoría de las revistas lo prohíben.
                  \item No es posible establecer un porcentaje fijo.
                  \item Algunas secciones son más críticas (resultados, conclusiones).
                  \item La IA no puede ser responsable.
                  \item Los humanos deben asumir responsabilidad completa por el contenido generado con IA.
                  \item ¿Dónde termina una simple corrección gramatical y comienza una reescritura semántica?.
              \end{itemize}

              %   \vspace{0.3cm}
              % \item \textbf{Consideraciones importantes:}
              %       \begin{itemize}
              %           \item Los revisores deben juzgar si la contribución es significativa.
              %           \item COPE debe fortalecer sus directrices sobre este tema.
              %           \item La contribución del autor debe ir más allá de editar texto generado por IA.
              %       \end{itemize}
    \end{itemize}

\end{frame}
\begin{frame}{La cuestión de la autoría}{\citeauthor{Lund2024}, \citeyear{Lund2024}}
    \citeauthor{Lund2024} en \cite{Lund2024} estudiaron las políticas de autoría de IA en \textit{300 journals}.
    % \begin{itemize}
    % \item 
    % \item De 300 revistas top analizadas, 58.7\% ya tenían políticas específicas sobre IA (en abril 2023).
    % \item Solo 3.4\% prohíben completamente el uso de IA (todas de la editorial Science)
    % \end{itemize}

    \begin{tikzpicture}
        \def\colores{cLabel_Dark,aLabel_Dark,aLabel_Dark,aLabel_Dark,aLabel_Dark}
        \def\words{300, 176,170,6,2}
        \def\texts{
            {revistas top analizadas},
            {ya tenían políticas específicas sobre IA (abril 2023)},
            {permiten usar IA para mejorar la calidad de manuscritos},
            {prohíben completamente el uso de IA (ed. Science)},
            {no prohíben explícitamente IA como autor (BioMed C.)}
        }

        \foreach [count=\index from 0] \word in \words{
            \ifnum\index=0
                \node[cLabel_Dark, text width=0.1\textwidth, align=center, minimum height=0.95cm] (text\index) at (0,-\index*1.1){\word};
            \else
                \node[aLabel_Dark, text width=0.1\textwidth, align=center, minimum height=0.95cm] (text\index) at (0,-\index*1.1){\word};
            \fi
        }
        \foreach [count=\index from 0] \text in \texts{
            \ifnum\index=0
                \node[cLabel, text width=0.78\textwidth, align=left, minimum height=1cm] (text\index)  at (5.3,-\index*1.1){\small \text};
            \else
                \node[aLabel, text width=0.78\textwidth, align=left, minimum height=1cm] (text\index)  at (5.3,-\index*1.1){\small \text};
            \fi
        }
    \end{tikzpicture}

\end{frame}
\begin{frame}{La cuestión de la autoría}{\citeauthor{Kendall2024}, \citeyear{Kendall2024}}
    \begin{itemize}
        \item Un autor puede declarar el uso de LLMs pero esto no dice cuánto lo ha usado.
        \item Tampoco que partes fueron hechas con la asistencia de LLMs.
        \item  Tanto el revisor como lector deben estar informados.
        \item \citeauthor{Kendall2024} \cite{Kendall2024} también hace una revisión de los \textit{guidelines} o directrices de seis revistas.
    \end{itemize}


\end{frame}
% Sección 2: Propuestas de Kendall
% \section{Propuestas de Kendall (2024)}
% \begin{frame}{Artículo de Kendall}{\fullcite{Kendall2024}}
%     \begin{itemize}
%         \item \textbf{Objetivo:} Propuestas para mejorar la transparencia en publicaciones científicas.
%         \item \textbf{Problema:} Directrices actuales no abordan el uso de IA, y son en algunos casos ambiguas.
%         \item \textbf{Propuesta:} Incluir detalles sobre el uso de IA en artículos.
%         \item \textbf{Importancia:} Mantener la integridad y confianza en la investigación científica.
%     \end{itemize}
%     \begin{boxA}{Cita clave}
%         \textit{“Sin transparencia, el fraude es difícil de detectar.”}
%     \end{boxA}
% \end{frame}
\section{Directrices de editoriales}
% Análisis de directrices de editoriales
\begin{frame}{Directrices de algunas editoriales}{En cuanto al uso de IA}
    \begin{threeparttable}
        \footnotesize
        \def\arraystretch{1.55}
        \begin{tabular}{lL{1.1cm}L{2.6cm}L{2.2cm}L{2.1cm}}
            \toprule
            {\footnotesize\textbf{Editorial}} & {\footnotesize Última versión} & \footnotesize\textbf{Reconocimiento obligatorio de uso} & \footnotesize\textbf{Especificar partes con IA}$^\red{1}$ & \footnotesize\textbf{Compartir prompts} \\
            \midrule
            BMJ$^\red{2}$ \cite{bmj_ai}       & 2024                           & \chekeYes                                               & \QuestionCircle                                           & \chekeYes                               \\
            Elsevier\cite{elsevier_ai}        & 2024                           & \QuestionCircle                                         & \QuestionCircle                                           & \chekeNot                               \\
            Sage$^\red{3}$\cite{sage_ai}      & 2023                           & \chekeYes                                               & \chekeYes                                                 & \chekeYes                               \\
            Springer\cite{springer_ai}        & 2024                           & \chekeYes                                               & \QuestionCircle                                           & \QuestionCircle                         \\
            Taylor \& F.\cite{taylor_ai}      & 2024                           & \chekeYes                                               & \chekeNot                                                 & \chekeNot                               \\
            Wiley\cite{wiley_ai}              & 2023                           & \chekeYes                                               & \chekeYes                                                 & \chekeYes                               \\
            \bottomrule
        \end{tabular}
    \end{threeparttable}
    \begin{tablenotes}\footnotesize\color{colorC}
        \item[] Directrices mencionadas por (\citeauthor{Kendall2024}, \citeyear{Kendall2024}) en \cite{Kendall2024}.
        \item[] $^\red{1}$ Muchas no lo especifican, pero requieren el \textit{``Why and how''}.
        \item[] $^\red{2}$ La más completa, según \cite{Kendall2024}
        \item[] $^\red{3}$ La más estricta, según \cite{Kendall2024}
        \item[] \QuestionCircle Recomienda o no especifica| \chekeYes Permite | \chekeNot No permite.
    \end{tablenotes}
\end{frame}
% Análisis de directrices de editoriales
\begin{frame}{Directrices de algunas editoriales}{En cuanto al uso de IA}
    \begin{threeparttable}
        \footnotesize
        \def\arraystretch{1.55}
        \begin{tabular}{lL{1.0cm}L{1.4cm}L{1.43cm}L{1.9cm}L{1.3cm}}
            % \begin{tabular}{cC{1.0cm}C{1.45cm}C{1.5cm}C{1.4cm}C{1.4cm}~}
            \toprule
            \textbf{Editorial}                     & \scriptsize\textbf{Última versión} & \scriptsize\textbf{Reconocer uso} & \scriptsize\textbf{Especificar partes} & \scriptsize\textbf{Compartir input/output} & \scriptsize\textbf{Imágenes y figuras$^\red{1}$} \\
            \midrule
            BMJ\red{$^2$} \cite{bmj_ai}            & 2024                               & \chekeYes                         & \QuestionCircle                        & \QuestionCircle                            & \QuestionCircle                                  \\
            Cambridge\red{$^2$}\cite{cambridge_ai} & 2023                               & \chekeYes                         & \chekeNot                              & \chekeNot                                  & \chekeYes                                        \\
            Elsevier     \cite{elsevier_ai}        & 2024                               & \chekeYes                         & \chekeNot                              & \chekeNot                                  & \chekeYes                                        \\
            IEEE\red{$^2$}    \cite{ieee_ai}       & 2023                               & \chekeYes                         & \chekeYes                              & \chekeYes                                  & \chekeYes                                        \\
            Oxford\red{$^2$} \cite{oxford_ai}      & 2023                               & \chekeYes                         & \QuestionCircle                        & \chekeNot                                  & \chekeYes                                        \\
            Sage            \cite{sage_ai}         & 2023                               & \chekeYes                         & \chekeYes                              & \chekeYes                                  & \chekeYes                                        \\
            Springer    \cite{springer_ai}         & 2024                               & \chekeYes                         & \QuestionCircle                        & \chekeNot                                  & \QuestionCircle                                  \\
            Taylor \& F. \cite{taylor_ai}          & 2024                               & \chekeYes                         & \QuestionCircle                        & \chekeNot                                  & \chekeNot                                        \\
            Wiley            \cite{wiley_ai}       & 2023                               & \chekeYes                         & \chekeYes                              & \chekeYes                                  & \chekeYes                                        \\
            \bottomrule
        \end{tabular}
    \end{threeparttable}
    \begin{tablenotes}\scriptsize\color{colorC}
        % \item[] Diferencias entre editoriales en sus políticas sobre el uso de IA.
        \item[]\red{1} Imágenes con carácter ilustrativo, no evidenciales.
        \item[]\red{2}Remite al \textit{Committee on Publication Ethics (COPE)}\cite{COPE_ai}.
        \item[] \QuestionCircle Recomienda o no especifica| \chekeYes Permite | \chekeNot No permite.
        % TF:  the full name of the tool used (with version number), how it was used, and the reason for use.
        %Only for improving language readability Springer, Elsevier
    \end{tablenotes}
\end{frame}
\begin{frame}{Directrices editoriales}{Declaración de uso de IA}
    \begin{itemize}
        \item Todas las editoriales permiten el uso de IA para mejorar la redacción (assistive AI).
        \item No todas especifican las partes generadas por IA.
    \end{itemize}
    \begin{tikzpicture}
        \def\words{Documented, How and Why, Prompts,Output, Autor}
        \foreach [count=\index from 0] \word in \words{
            \node[aLabel_Dark, text width=0.25\textwidth, align=left] (text\index)  at (0,-\index*0.95){\word};
        }
        \def\words{Springer-Wiley-Cambridge-IEEE-Sage, Taylor \& F. - Elsevier - Springer - Wiley, Sage - IEEE,Sage - IEEE, Ninguna}
        \foreach [count=\index from 0] \word in \words{
            \node[aLabel, text width=0.60\textwidth, align=left] (text\index)  at (5.15,-\index*0.95){\word};
        }
        \node[dLabel, text width=0.9\textwidth, align=left] (text\index)  (Bottom) at (3.5,-5){Sage Publishers provee un template obligatorio para autores que usan Generative AI \cite{sage_ai}.};
    \end{tikzpicture}


\end{frame}
\begin{frame}{Directrices editoriales}{Distinción entre el carácter asistencial y la IA generativa}
    \begin{itemize}
        \item Cada vez más editoriales distinguen explícitamente entre el uso de IA asistencial (Assistive AI) y IA generativa (Generative AI).
        \item Es importante declarar el uso asistido de LLM's
    \end{itemize}
    \begin{tikzpicture}
        \node[cLabel_Dark, text width=0.41\textwidth, align=left] (base) at (0,0) {Assistive AI };
        \def\words{Corrección de texto, Mejora de redacción, Revisión de estilo, Crear diagramas, No declarar}
        \foreach [count=\index from 1] \word in \words{
            \node[cLabel, text width=0.41\textwidth, align=left] (text\index)  at (0,-\index*1.0){{\color{GithubGreen}\faCheckCircle} \word};
        }
        \node[dLabel_Dark, text width=0.41\textwidth, align=left] (text0) at (6,0) {Generative AI };
        \def\words{Crear contenido, Generar texto, Crear imágenes, Generar gráficos, Obligatorio declarar}
        \foreach [count=\index from 1] \word in \words{
            \node[dLabel, text width=0.41\textwidth, align=left] (text\index)  at (6,-\index*1.0){{\color{colorD}\faTimesCircle} \word};
        }
    \end{tikzpicture}
\end{frame}

\begin{frame}{Directrices Propuestas}
    \begin{enumerate}
        \item \alerta{Declaración detallada:} Herramienta, versión y fechas de uso.
        \item \alerta{Especificar las secciones fueron generadas por IA.}
        \item \alerta{Incluir prompts y texto generado} en archivo suplementario.
        \item \alerta{Incluir texto original} si se usó IA para editar.
        \item \alerta{Archivo suplementario} con toda la información.
        \item \alerta{Revisión por pares informada:} Comparar texto original vs. editado.
        \item \alerta{Hacerlo accesible} a revisores y lectores.
        \item \alerta{Directrices unificadas} (ej.: COPE + políticas editoriales).
    \end{enumerate}
    \begin{boxC}{Importante}
        Se trata de \textbf{restaurar la confianza} en el proceso científico.
    \end{boxC}
\end{frame}

% Propuesta: Nuevas directrices (1)
\begin{frame}{Propuestas}{Nuevas directrices}
    \begin{tikzpicture}
        \def\words{
            Reconocimiento completo,
            Herramienta utilizada,
            Versión específica,
            Fechas de uso}

        \foreach [count=\index from 0] \word in \words{
            \ifnum\index=0
                \node[aLabel_Dark, text width=0.45\textwidth, align=center, minimum height=0.6cm] (text\index) at (0,-\index*0.9){\word};
            \else
                \node[aLabel, text width=0.45\textwidth, align=center, minimum height=0.6cm] (text\index) at (0,-\index*0.9){\word};
            \fi
        }
        \begin{scope}[xshift=5.65cm]
            \def\words{
                Especificar secciones,
                Partes generadas con IA
            }

            \foreach [count=\index from 0] \word in \words{
                \ifnum\index=0
                    \node[aLabel_Dark, text width=0.45\textwidth, align=center, minimum height=0.6cm] (text\index) at (0,-\index*0.9){\word};
                \else
                    \node[aLabel, text width=0.45\textwidth, align=center, minimum height=0.6cm] (text\index) at (0,-\index*0.9){\word};
                \fi
            }
        \end{scope}
        \begin{scope}[yshift=-3.85cm]
            \def\words{
                Transparencia en los prompts,
                Incluir prompts utilizados,
                Mostrar texto generado
            }

            \foreach [count=\index from 0] \word in \words{
                \ifnum\index=0
                    \node[cLabel_Dark, text width=0.48\textwidth, align=center, minimum height=0.6cm] (text\index) at (0,-\index*0.9){\word};
                \else
                    \node[cLabel, text width=0.48\textwidth, align=center, minimum height=0.6cm] (text\index) at (0,-\index*0.9){\word};
                \fi
            }
        \end{scope}
        % \item \textbf{Mejoras de lenguaje}:
        % \begin{itemize}
        %     \item Reconocer si la IA se utilizó para mejorar la redacción
        %     \item Mostrar el texto original antes de la mejora
        % \end{itemize}
    \end{tikzpicture}
\end{frame}

% Propuesta: Nuevas directrices (2)
\begin{frame}{Propuestas}{Nuevas directrices}
    \begin{tikzpicture}
        \def\words{
            Mejoras de lenguaje,
            Reconocer se utilizó para mejorar redacción,
            Mostrar el texto antes de la mejora
        }

        \foreach [count=\index from 0] \word in \words{
            \ifnum\index=0
                \node[aLabel_Dark, text width=0.8\textwidth, align=center, minimum height=0.45cm] (text\index) at (0,-\index*0.3){\word};
            \else
                \node[aLabel, text width=0.8\textwidth, align=left, minimum height=0.45cm] (text\index) at (0,-\index*1.1){\small \word};
            \fi
        }

        \begin{scope}[yshift=-3.75cm]
            \def\words{
                Archivo suplementario,
                Incluir todos los detalles anteriores,
                El archivo debe estar disponible para revisores y lectores,
                Permite evaluación informada del aporte real
            }

            \foreach [count=\index from 0] \word in \words{
                \ifnum\index=0
                    \node[cLabel_Dark, text width=0.8\textwidth, align=center, minimum height=0.45cm] (text\index) at (0,-\index*0.1){\small \word};
                \else
                    \node[cLabel, text width=0.8\textwidth, align=left, minimum height=0.45cm] (text\index) at (0,-\index*0.9){\small \word};
                \fi
            }
        \end{scope}
    \end{tikzpicture}
\end{frame}

% Beneficios de la propuesta
\begin{frame}{Beneficios de la propuesta}
    \begin{itemize}
        \item \textbf{Para revisores}:
              \begin{itemize}
                  \item Evaluación informada sobre el aporte real del autor
                  \item Mejor criterio para juzgar la autoría
              \end{itemize}
              \vspace{0.3cm}
        \item \textbf{Para lectores}:
              \begin{itemize}
                  \item Transparencia sobre el origen del contenido
                  \item Confianza en la integridad de la publicación
              \end{itemize}
              \vspace{0.3cm}
        \item \textbf{Para la ciencia}:
              \begin{itemize}
                  \item Integridad en la autoría y en el proceso de publicación
                  \item Estándares más claros y consistentes
              \end{itemize}
    \end{itemize}
\end{frame}
% Sección 3: Casos de Estudio
\section{Casos de Estudio}
\begin{frame}{Casos de Estudio}
    \begin{tikzpicture}
        \def\words{
            Som Biswas,
            {\href{https://orcid.org/0000-0002-4038-5844}{\color{GithubGreen}\faOrcid \hspace{1em}{\small https://orcid.org/0000-0002-4038-5844}}},
            {Publica 10 artículos en 2023 donde declara haber usado inteligencia artificial, pero no dice cuánto.},
            {\textit{``According to Chat GPT''}}
        }

        \foreach [count=\index from 0] \word in \words{
            \ifnum\index=0
                \node[aLabel_Dark, text width=0.8\textwidth, align=left, minimum height=0.45cm] (text\index) at (0,-\index*0.3){\word};
            \else
                \node[aLabel, text width=0.8\textwidth, align=left, minimum height=0.45cm] (text\index) at (0,-\index*1.1){\small \word};
            \fi
        }

        \begin{scope}[yshift=-5.05cm]
            \def\words{
                {M. Majovský},
                Elabora un artículo médico fraudulento con ChatGPT que engaña a expertos
            }
            \foreach [count=\index from 0] \word in \words{
                \ifnum\index=0
                    \node[cLabel_Dark, text width=0.8\textwidth, align=left, minimum height=0.45cm] (text\index) at (0,-\index*0.1){\small \word};
                \else
                    \node[cLabel, text width=0.8\textwidth, align=left, minimum height=0.45cm] (text\index) at (0,-\index*1.1){\small \word};
                \fi
            }
        \end{scope}
    \end{tikzpicture}
\end{frame}
\begin{frame}{Caso 1}{S. Biswas (2023)}
    \begin{itemize}
        \item \textbf{Experimento:} Generación de artículo médico fraudulento con ChatGPT.

        \item Publicó 13 artículos con ``ChatGPT'' en el título durante 2023
        \item Reconoció, sin especificar lo siguiente:
              \begin{center}
                  \begin{tikzpicture}
                      \node[aLabel, text width=0.8\textwidth,align=left] (a) at (0,0) {\textit{``Este artículo fue parcialmente generado por ChatGPT... La edición fue realizada por el autor''}};
                  \end{tikzpicture}
              \end{center}
        \item \textbf{Problemas identificados}:
              \begin{itemize}
                  \item No sabemos qué porcentaje del texto fue generado por IA
                  \item No sabemos qué ediciones hizo el autor
                  \item Los artículos cubren temas muy diversos (inusual para un solo autor)
                  \item No cita sus propios trabajos (práctica común en publicaciones académicas)
              \end{itemize}
    \end{itemize}


\end{frame}

\begin{frame}[allowframebreaks]{Caso 1}{S. Biswas (2023)}
    \begin{threeparttable}
        % \centering
        \scriptsize
        \def\arraystretch{1.75}
        % \caption{\footnotesize Listados de artículos por S. Biswas}
        \vspace{-1em}
        \begin{tabular}{c|L{8.75cm}|C{0.55cm}} \toprule
            \textbf{Año}            & \textbf{Título del artículo} & \textbf{Cita}       \\\hline
            \citeyear{biswas_2023d} & \citetitle{biswas_2023d}     & \cite{biswas_2023d} \\
            \citeyear{biswas_2023e} & \citetitle{biswas_2023e}     & \cite{biswas_2023e} \\
            \citeyear{biswas_2023j} & \citetitle{biswas_2023j}     & \cite{biswas_2023j} \\
            \citeyear{biswas_2023k} & \citetitle{biswas_2023k}     & \cite{biswas_2023k} \\
            \citeyear{biswas_2023a} & \citetitle{biswas_2023a}     & \cite{biswas_2023a} \\
            \citeyear{biswas_2023b} & \citetitle{biswas_2023b}     & \cite{biswas_2023b} \\
            \citeyear{biswas_2023i} & \citetitle{biswas_2023i}     & \cite{biswas_2023i} \\
            \citeyear{biswas_2023c} & \citetitle{biswas_2023c}     & \cite{biswas_2023c} \\
            \citeyear{biswas_2023f} & \citetitle{biswas_2023f}     & \cite{biswas_2023f} \\
            \citeyear{biswas_2023h} & \citetitle{biswas_2023h}     & \cite{biswas_2023h} \\
            \bottomrule
        \end{tabular}
        \vspace{0.25em}
        \begin{tablenotes}
            \item[] {\color{colorD}Todos los artículos fueron publicados en 2023 y el autor reconoce la asistencia de ChatGPT.}
        \end{tablenotes}
    \end{threeparttable}
\end{frame}

\begin{frame}{Caso 2}{Májovský et al. (2023)}

    \begin{itemize}
        \item \textit{Proof-of-concept study}\cite{majovsky2023}:
              El estudio logró generar un artículo completo (introducción, métodos, resultados, discusión, referencias y gráficos) con ChatGPT (GPT-3).
        \item Usando 13 prompts, < 1 h de trabajo humano.
        \item \alerta{Hallazgos:}
              \begin{itemize}
                  \item Estructura impecable (abstract, métodos, resultados).
                  \item Citas falsas detectadas por expertos.
                  \item Datos inconsistentes en tablas estadísticas.
                  \item Ausencia de registros en ClinicalTrials.gov, debilidades metodológicas.
              \end{itemize}
        \item \alerta{Revisores:} Sometido a detectores de IA, evaluación experta, estadísticos y ChatGPT.
        \item \alerta{Citas:} Hizo \alert{17 citas}, 9 de ellas correctas, 4 con errores menores y \alert{4 inexistentes}.
        \item \alerta{Conclusión:} Sin transparencia, el fraude es difícil de detectar.
    \end{itemize}
    % \begin{center}
    % \end{center}
    % \begin{tikzpicture}
    %     \node[aLabel, text width=0.99\textwidth,align=left] (a) at (0,0) {\footnotesize\fullcite{majovsky2023}};
    % \end{tikzpicture}
\end{frame}

\begin{frame}{Otros ejemplos - \citeauthor{walters_fabrication_2023},\citeyear{walters_fabrication_2023}}

    \begin{tikzpicture}
        \def\colores{cLabel_Dark,aLabel_Dark,aLabel_Dark,aLabel_Dark,aLabel_Dark}
        \def\words{Enfoque, 42,636,84}
        \def\texts{
            {Errores en citas bibliográficas falsas o incorrectas generadas por ChatGPT-3.5 y ChatGPT-4\cite{walters_fabrication_2023} \alerta{(haciendo lit. reviews)}.},
            {Temas multidisciplinarios},
            {citas encontradas},
            {artículos en total}
        }

        \foreach [count=\index from 0] \word in \words{
            \ifnum\index=0
                \node[cLabel_Dark, text width=0.15\textwidth, align=center, minimum height=0.85cm] (text\index) at (0,-\index*1.1){\word};
            \else
                \node[aLabel_Dark, text width=0.15\textwidth, align=center, minimum height=0.85cm] (text\index) at (0,-\index*1.1){\word};
            \fi
        }
        \foreach [count=\index from 0] \text in \texts{
            \ifnum\index=0
                \node[cLabel, text width=0.8\textwidth, align=left, minimum height=0.85cm] (text\index)  at (5.4,-\index*1.1){\footnotesize \text};
            \else
                \node[aLabel, text width=0.8\textwidth, align=left, minimum height=0.85cm] (text\index)  at (5.4,-\index*1.1){\footnotesize \text};
            \fi
        }
    \end{tikzpicture}


    \begin{columns}
        \column{0.65\textwidth}
        \begin{itemize}
            \item \alerta{ChatGPT-3.5:}
                  \begin{itemize}
                      \item 55\% citas fabricadas.
                      \item 43\% con errores sustanciales.
                  \end{itemize}
            \item \alerta{ChatGPT-4:}
                  \begin{itemize}
                      \item 18\% citas fabricadas.
                      \item 24\% con errores.
                  \end{itemize}
        \end{itemize}
        \column{0.5\textwidth}
        \begin{table}[ht]
            \centering
            \footnotesize
            \begin{tabular}{@{}ll@{}}
                \toprule
                \textbf{Error}      & \textbf{Frecuencia} \\
                \midrule
                Autores incorrectos & 14\%                \\
                Títulos falsos      & 6\%                 \\
                Fechas erróneas     & 22\%                \\
                \bottomrule
            \end{tabular}
        \end{table}
    \end{columns}
\end{frame}

\begin{frame}{Errores de Citación}{\citeauthor{walters_fabrication_2023},\citeyear{walters_fabrication_2023}}
    \resizebox{0.85\textwidth}{!}{%
        \begin{tikzpicture}
            \begin{axis}[
                    ybar,
                    symbolic x coords={
                            Autores,
                            Títulos,
                            Fechas,
                            Journals,
                            Vol/Pags,
                            Publisher
                        },
                    xtick=data,
                    ymin=0, ymax=38,
                    ylabel={Porcentaje (\%)},
                    nodes near coords,
                    bar width=10pt,
                    enlarge x limits=0.1,
                    legend style={
                            at={(0.75,1.15)},
                            anchor=north,
                            legend columns=-1
                        },
                    xticklabel style={rotate=45, anchor=east},
                    tick label style={font=\small},
                    label style={font=\scriptsize},
                ]
                \addplot coordinates {
                        (Autores, 14)
                        (Títulos, 6)
                        (Fechas, 22)
                        (Journals, 14)
                        (Vol/Pags, 34)
                        (Publisher, 9)
                    };
                \addplot coordinates {
                        (Autores, 6)
                        (Títulos, 3)
                        (Fechas, 16)
                        (Journals, 4)
                        (Vol/Pags, 13)
                        (Publisher, 3)
                    };
                \legend{GPT-3.5, GPT-4}

            \end{axis}


            \node[dLabel, rotate=90, xshift=35mm] (ref) at (7.3, -2){\scriptsize Con los datos de \cite{walters_fabrication_2023}};
        \end{tikzpicture}%
    }
\end{frame}
% Sección 4: Recomendaciones
\section{Recomendaciones}
\begin{frame}{Implementando las Propuestas}
    \begin{boxA}{Para Autores}
        \begin{itemize}
            \item Adjuntar \alerta{archivo suplementario} con prompts y texto generado.
            \item Especificar secciones editadas (ej.: métodos, discusión).
        \end{itemize}
    \end{boxA}

    \begin{boxA}{Para Revistas}
        \begin{itemize}
            \item Adoptar políticas como \alerta{Wiley} o \alerta{Sage}:
                  \begin{itemize}
                      \item Plantillas para declarar uso de IA.
                      \item Verificación manual de citas clave.
                  \end{itemize}
            \item Capacitar revisores en detección de IA.
        \end{itemize}
    \end{boxA}
\end{frame}

% Sección 5: Conclusión
\section{Conclusión}
\begin{frame}{Conclusión}
    \begin{itemize}
        \item \alerta{Riesgos:} Fraude, pérdida de confianza, citas falsas.
        \item \alerta{Oportunidades:} IA mejora eficiencia si se usa con transparencia.
        \item \alerta{Llamado a la acción:}
              \begin{itemize}
                  \item Adoptar directrices unificadas \alerta{y explícitas}.
                  \item[\color{colorD}\faArrowRight] ej.: %COPE + ICMJE + \citeauthor{Kendall2024},\citeyear{Kendall2024}

                        \begin{center}
                            \begin{tikzpicture}
                                \node[dLabel, text width=0.1\textwidth, align=left] (origin)  at (0,0){\small\bfseries COPE};
                                \node[text=colorA,font=\bfseries,right=1mm of origin] (a)  {+};
                                \node[dLabel, text width=0.1\textwidth, align=left, right=1mm of a] (b)  {\small\bfseries ICMJE};
                                \node[text=colorA,font=\bfseries,right=1mm of b] (c)  {+};
                                \node[cLabel, text width=0.2\textwidth, align=left, right=1mm of c] (d)  {\small \citeauthor{Kendall2024}, \citeyear{Kendall2024}};
                            \end{tikzpicture}
                        \end{center}
                  \item Fomentar cultura de responsabilidad académica.
              \end{itemize}
    \end{itemize}
    \centering
    % \includegraphics[width=0.3\textwidth]{transparency_icon.png} % Añade una imagen relevante
\end{frame}

% Conclusiones
\begin{frame}{Conclusiones}
    \begin{itemize}
        \item Las directrices actuales son inconsistentes e insuficientes
        \item Se necesita transparencia total sobre:
              \begin{itemize}
                  \item Qué partes fueron generadas con IA
                  \item Los prompts utilizados
                  \item El texto generado antes de la edición
              \end{itemize}
              \vspace{0.3cm}
        \item La comunidad científica debe adoptar estándares más robustos
        \item El objetivo: mantener la integridad y la confianza en las publicaciones científicas
    \end{itemize}
\end{frame}
\begin{frame}[allowframebreaks]{Referencias}{}
    \scriptsize
    % \printbibliography[heading=none]
    \printbibliography[keyword={ethics},title={Paginas web}]
\end{frame}

\begin{frame}[allowframebreaks]{Referencias}{Directrices de editoriales}
    \tiny
    % \printbibliography[heading=none]
    \printbibliography[keyword={online},title={Paginas web}]
\end{frame}


\setbeamertemplate{logo}{}
\begin{frame}[allowframebreaks]
    % {Artículos de Son Biswas de 2023}{Generados parcial o totalmente con ChatGPT}
    \tiny
    \defbibnote{nota3}{Artículos de S. Biswas, generados parcial o totalmente con ChatGPT:}

    \printbibliography[title={}, prenote=nota3, keyword={example}, heading=subbibliography]


\end{frame}
% Referencias



\begin{frame}
    \begin{enumerate}\small
        \item \fullcite{Kendall2024}
        \item \fullcite{walters_fabrication_2023}
        \item \fullcite{majovsky2023}
    \end{enumerate}
\end{frame}

% \begin{frame}
%     \frametitle{Introduction: The Rise of AI in Research}
%     \begin{itemize}
%         \item Use of Large Language Models (LLMs) like ChatGPT in academic writing.
%         \item Tools like ChatGPT are now used for:
%               \begin{itemize}
%                   \item Drafting papers
%                   \item Editing language
%                   \item Generating ideas
%               \end{itemize}
%         \item Central question: How do we ensure transparency and integrity when AI is involved in writing scientific papers?
%     \end{itemize}
% \end{frame}

% \begin{frame}
%     \frametitle{Kendall's Core Argument}
%     \begin{itemize}
%         \item Current journal guidelines are not sufficient to ensure transparency in AI-assisted writing.
%         \item Key proposal:
%               \begin{itemize}
%                   \item Authors should disclose prompts, generated text, and edits.
%                   \item This should be available to reviewers and readers.
%               \end{itemize}
%         \item Ethical concern: A paper could be 100\% AI-generated and still meet current disclosure standards.
%     \end{itemize}
% \end{frame}

% \begin{frame}
%     \frametitle{Case Study - Biswas}
%     \begin{itemize}
%         \item S. Biswas published over a dozen papers in 2023 with ChatGPT assistance.
%         \item Issues:
%               \begin{itemize}
%                   \item Vague acknowledgments like “partially generated by ChatGPT.”
%                   \item No clarity on which parts were AI-generated or how much editing was done.
%                   \item Topics ranged widely (from gaming to public health), raising questions about authorship and originality.
%               \end{itemize}
%         \item Kendall's critique: This is exactly the kind of case where stronger guidelines are needed.
%     \end{itemize}
% \end{frame}

% \begin{frame}
%     \frametitle{Case Study - Majovský}
%     \begin{itemize}
%         \item Majovský et al. (2023) experiment:
%               \begin{itemize}
%                   \item Used ChatGPT to generate a completely fake medical article in 1 hour.
%                   \item Included fabricated data, references, and even a chart.
%                   \item Experts found it convincing, but flawed on closer inspection.
%               \end{itemize}
%         \item Key takeaway: AI can produce authentic-looking fraud. Detection is hard. Reviewers need more information.
%     \end{itemize}

% \end{frame}

% \begin{frame}
%     \frametitle{Walters and Wilder: Citation Fabrication}
%     \begin{itemize}
%         \item Findings from Walters and Wilder (2023):
%               \begin{itemize}
%                   \item GPT-3.5: 55\% of citations were fake.
%                   \item GPT-4: Improved, but still 18\% fabricated.
%                   \item Even real citations had errors in 24-43\% of cases.
%               \end{itemize}
%         \item Implication: AI-generated references are unreliable. Without transparency, readers and reviewers are misled.
%     \end{itemize}
% \end{frame}

% \begin{frame}
%     \frametitle{Conclusion - A Call to Action}
%     \begin{itemize}
%         \item Reiterate the stakes: AI is here to stay. But without transparency, we risk eroding the credibility of science.
%         \item End with a powerful quote or question: “Should we trust what we read, if we don't know who—or what—wrote it?”
%     \end{itemize}
% \end{frame}


% \begin{frame}{Percentage of Fabricated Citations}
%     \begin{tikzpicture}
%         \begin{axis}[
%                 ybar,
%                 symbolic x coords={GPT-3.5, GPT-4},
%                 xtick=data,
%                 ymin=0, ymax=60,
%                 ylabel={Percentage (\%)},
%                 nodes near coords,
%                 bar width=20pt,
%                 enlarge x limits=0.5
%             ]
%             \addplot coordinates {(GPT-3.5, 55) (GPT-4, 18)};
%         \end{axis}
%     \end{tikzpicture}
% \end{frame}

% \begin{frame}{Percentage of Citation Errors Among Real Citations}
%     \begin{tikzpicture}
%         \begin{axis}[
%                 ybar,
%                 symbolic x coords={GPT-3.5, GPT-4},
%                 xtick=data,
%                 ymin=0, ymax=50,
%                 ylabel={Percentage (\%)},
%                 nodes near coords,
%                 bar width=20pt,
%                 enlarge x limits=0.5
%             ]
%             \addplot coordinates {(GPT-3.5, 43) (GPT-4, 24)};
%         \end{axis}
%     \end{tikzpicture}
% \end{frame}

% \begin{frame}{Types of Citation Errors}
%     \begin{tikzpicture}
%         \begin{axis}[
%                 ybar,
%                 symbolic x coords={Incorrect Author Names, Incorrect Titles, Incorrect Dates, Incorrect Journal Titles, Incorrect Volume/Page Numbers, Incorrect Publisher Names},
%                 xtick=data,
%                 ymin=0, ymax=40,
%                 ylabel={Percentage (\%)},
%                 nodes near coords,
%                 bar width=10pt,
%                 enlarge x limits=0.2,
%                 legend style={at={(0.5,0.15)}, anchor=north, legend columns=-1},
%                 xticklabel style={rotate=45, anchor=east}
%             ]
%             \addplot coordinates {(Incorrect Author Names, 14) (Incorrect Titles, 6) (Incorrect Dates, 22) (Incorrect Journal Titles, 14) (Incorrect Volume/Page Numbers, 34) (Incorrect Publisher Names, 9)};
%             \addplot coordinates {(Incorrect Author Names, 6) (Incorrect Titles, 3) (Incorrect Dates, 16) (Incorrect Journal Titles, 4) (Incorrect Volume/Page Numbers, 13) (Incorrect Publisher Names, 3)};
%             \legend{GPT-3.5, GPT-4}
%         \end{axis}
%     \end{tikzpicture}
% \end{frame}

